\section{Information-Theoretic Interpretation of Internal State}
\label{sec:information_theory}

\paragraph{Model Taxonomy.}
We use two classes, both with the same locked endpoint readout
\[
\Delta P = k\,[m(\tau_T)-m(\tau_0)].
\]
\begin{itemize}
\item \textbf{Class I (Invariant leaky accumulator / Paper-1):}
\[
\frac{dm}{d\tau}=\beta-\lambda m,\qquad
\Delta P = k\,[m(\tau_T)-m(\tau_0)].
\]
The drive is $\tau$-intrinsic.
\item \textbf{Class II (Coordinate-surprise observer):}
\[
\frac{dm}{d\tau}=\beta I(\tau)-\lambda m,\qquad
I(\tau)=\left|\frac{d\log(dt/d\tau)}{d\tau}\right|,\qquad
\Delta P = k\,[m(\tau_T)-m(\tau_0)].
\]
Only the drive differs from Class I.
\end{itemize}

\paragraph{Proposition (Generic coordinate dependence of $I(\tau)$).}
Let $t\mapsto s(t)$ be a strictly monotone reparameterization of coordinate time. Then $r(\tau)$ transforms as
\begin{equation}
r(\tau)\mapsto r_s(\tau)=\frac{ds}{d\tau}=\frac{ds}{dt}\frac{dt}{d\tau}=\frac{ds}{dt}\,r(\tau).
\label{eq:r_transform}
\end{equation}
Hence
\begin{equation}
\frac{d\log r_s}{d\tau}
=
\frac{d\log r}{d\tau}
+
\frac{d}{d\tau}\log\!\left(\frac{ds}{dt}\right),
\label{eq:logr_transform}
\end{equation}
so $I(\tau)=|d\log r/d\tau|$ is not invariant in general (except special cases such as affine $s$).

\paragraph{Proof sketch.}
Equation~\eqref{eq:r_transform} is the chain rule. Taking $\tau$-derivatives of logarithms gives Eq.~\eqref{eq:logr_transform}. The extra additive term is generically nonzero for nonlinear $s(t)$, therefore $I(\tau)$ and any model driven by it are coordinate dependent.
Exception: if $r(\tau)$ is constant along the trajectory family, then $I(\tau)\equiv 0$ and the coordinate-dependent drive term vanishes.

\paragraph{Consequence for matched-endpoint comparisons.}
For matched endpoints $(\tau_0,\tau_T)$, Class I remains invariant up to numerical precision in warp tests, while Class II and fixed-$t$ observers (H1/H2) exhibit finite gauge violation.
Gauge invariance is guaranteed only for $\tau$-intrinsic drives; coordinate-time drives may be invariant on special families ($r$ constant) but are not invariant generically.
